
\documentclass[aspectratio=169]{beamer}

\mode<presentation>
\usetheme{Boadilla}
\definecolor{NTHU1}{RGB}{127,16,132}
\definecolor{NTHU2}{RGB}{228,210,231}
\definecolor{blue}{RGB}{30,90,205}
\definecolor{red}{RGB}{213,94,0}
\definecolor{green}{RGB}{0,128,0}
\definecolor{charcoalgray}{RGB}{108,104,104}
\setbeamercolor{title}{fg=NTHU1}
\setbeamercolor{frametitle}{fg=NTHU1}
\setbeamercolor{block title}{bg=NTHU1, fg=white}
\setbeamercolor{block body}{bg=NTHU2}
\setbeamercolor{structure}{fg=NTHU1}
\setbeamercolor{item projected}{fg=white}
\setbeamercolor{item}{fg=NTHU1}
\setbeamercolor{subitem}{fg=NTHU1}
\setbeamercolor{section in toc}{fg=NTHU1}
\setbeamercolor{description item}{fg=NTHU1}
\setbeamercolor{caption name}{fg=NTHU1}
\setbeamercolor{button}{bg=NTHU1, fg=white}
\setbeamercolor{caption name}{fg=NTHU1}
\usepackage{graphics}
\usepackage{tikz}
\usepackage{amsmath}
\usepackage{bbm}
\usetikzlibrary{decorations.pathreplacing}
\usepackage{geometry}
\usepackage{booktabs}
\usepackage{bookmark}
\usepackage{multirow, makecell}
\usepackage{float}
\usepackage{fancyvrb}
\usepackage{caption}
\captionsetup[figure]{singlelinecheck = false, format= hang, justification=centering}
\captionsetup[subfigure]{singlelinecheck = false, format= hang, justification=raggedright}
\usepackage{subcaption}
\usepackage{adjustbox}
\usepackage{hyperref}
\usepackage{threeparttable}
\usepackage{appendixnumberbeamer} 
\usepackage[T1]{fontenc}
\usepackage[default]{lato} 
\usepackage{smartdiagram}
\smartdiagramset{
  back arrow disabled = true,
  module minimum width=1cm,
  module x sep = 3,
  uniform arrow color=true,
}
\usepackage{xcolor}
\usepackage{colortbl}
  \definecolor{light-gray}{gray}{0.99}

\newenvironment{wideitemize}{\itemize\addtolength{\itemsep}{10pt}}{\enditemize}
\newenvironment{wideenumerate}{\enumerate\addtolength{\itemsep}{10pt}}{\endenumerate}
\newenvironment{widedescription}{\description\addtolength{\itemsep}{10pt}}{\enddescription}
\hypersetup{
colorlinks=true,
linkcolor=NTHU1,
filecolor=green, 
urlcolor=blue,
}
\beamertemplatenavigationsymbolsempty
\setbeamercolor{author in head/foot}{bg=white, fg=NTHU1}
\setbeamercolor{title in head/foot}{bg=white, fg=NTHU1}
\setbeamercolor{date in head/foot}{bg=white, fg=NTHU1}
\setbeamercolor{section in head/foot}{bg=white, fg=NTHU1}
\setbeamercolor{headline}{bg=white}
\setbeamertemplate{footline}{
    \leavevmode%
    \hbox{%
        \begin{beamercolorbox}[wd=.333333\paperwidth,ht=2.25ex,dp=1ex,center]{date in head/foot}%
            \usebeamerfont{date in head/foot}%\insertdate
        \end{beamercolorbox}%
        \begin{beamercolorbox}[wd=.333333\paperwidth,ht=2.25ex,dp=1ex,center]{title in head/foot}%
            \usebeamerfont{title in head/foot}\insertshorttitle
        \end{beamercolorbox}%
        \begin{beamercolorbox}[wd=.333333\paperwidth,ht=2.25ex,dp=1ex,right]{date in head/foot}%
            \usebeamerfont{date in head/foot}\insertframenumber{}\hspace*{2em}
        \end{beamercolorbox}}%
        \vskip0pt%
    }
    %% May need to script the introduction!!!!!!
%\setbeamercolor{page number in head/foot}{fg=NTHU1}
\setbeamertemplate{section in toc}[sections numbered]
\setbeamertemplate{subsection in toc}{\leavevmode\leftskip=3em\rlap{\hskip-1.75em\inserttocsectionnumber.\inserttocsubsectionnumber}
\inserttocsubsection\par}
\setbeamerfont{subsection in toc}{size=\footnotesize}
%\setbeamertemplate{headline}{%
  %\begin{beamercolorbox}[ht=5.5ex]{section in head/foot}
    %\vskip2pt\insertnavigation{0.33\paperwidth}\vskip2pt
  %\end{beamercolorbox}%
%}
\newenvironment{transitionframe}{\setbeamercolor{background canvas}{bg=white}\setbeamertemplate{footline}{} \begin{frame}[noframenumbering]}{\end{frame}}

\makeatletter
\let\@@magyar@captionfix\relax
\makeatother

\title[]{Conflict and Peace} % Change this regularly
\subtitle[]{Week 4: Engaging in (organized) crime}
\author[Lee]{Seung-hun Lee }
\institute[]{Taipei School of Economics\\ National Tsing Hua University}

\date[]{September 26th, 2025}

\begin{document}
\begin{transitionframe}
\titlepage
\end{transitionframe}

%%% Color slides for section headers: Use for colloquium version (The ones bewteen \iffals and \fi)

\section{Overview of the topic}
\begin{frame}
\frametitle{What drives people into crime?} 
We have already highlighted how economic forces motivate people into conflicts
\begin{itemize}\setlength\itemsep{3pt}
\item If opportunity cost of work is higher, people less likely to join crime
\begin{itemize}
\item Employment prospects, Returns from education, long-run income flow discourage crime
\end{itemize}
\item People often compare their returns in legal markets vs illegal markets
 \begin{itemize}
\item This explains many types of crimes: From joining criminal groups to tax evasion
\end{itemize}
\end{itemize}\medskip
Here we focus on organized criminal groups
\begin{itemize}\setlength\itemsep{3pt}
\item Not just the violence they cause, but also the strategic motives behind their business
\item Also have recruiting mechanisms that target the vulnerable
 \begin{itemize}
\item Youth and the poor, sometimes (ethnically/culturally) marginalized individuals
\end{itemize}
\item Sometimes run their own `governments' (criminal governance)
\begin{itemize}
\item Collect ``protection fees" (extortions), compete with formal governments for citizen support
\end{itemize}
\end{itemize}\medskip
\end{frame}

\begin{frame}
\frametitle{Youth disproportionately affected} 
\begin{figure}[H]
  \includegraphics[keepaspectratio,width=0.6\textwidth]{fig3.png}
    \captionsetup{justification=raggedright,singlelinecheck=false}
\caption*{\footnotesize{Source: Soares, Rodrigo (2006) ``The welfare cost of violence across countries", \textit{Journal of Health Economics}, 25, 821-846}}
\end{figure}

\end{frame}

\begin{frame}
\frametitle{Large share of criminal groups, but with short career paths} 
\begin{figure}[H]
  \includegraphics[keepaspectratio,width=0.6\textwidth]{fig2.png}
    \captionsetup{justification=raggedright,singlelinecheck=false}
\caption*{\footnotesize{Source: Preito-Curiel et al (2023) ``Reducing cartel recruitment is the only ay to lower violence in Mexico", \textit{Science}, 381, 1312-1316 }}
\end{figure}
\end{frame}

\begin{frame}
\frametitle{Criminal governance is actually widespread} 
What is criminal governance?
\begin{itemize}\setlength\itemsep{3pt}
\item Provision of order, services, and market regulation by non-state criminal entities
\item Often in places with absent or captured (central/local) governments 
\begin{itemize}\setlength\itemsep{3pt}
\item Bandits collecting mineral taxation in Congo (Sanchez de la Sierra 2020)
\item Italian Mafia responsible for trash collection (some estimates say \$ 8.8B USD worth)
\end{itemize}
\item Replaces state in some cases: Difficult to hold criminals accountable
\end{itemize}\medskip
Continuously contested question: Why?
\begin{itemize}\setlength\itemsep{3pt}
\item Why are some states unable to address criminal governance?
\item Why do organized criminal groups engage in governance, not direct violence (or mix)?
\end{itemize}
\end{frame}

\begin{frame}
\frametitle{Criminal governance is actually widespread} 
\begin{figure}[H]
  \includegraphics[keepaspectratio,width=0.67\textwidth]{fig4.png}

\end{figure}
\end{frame}

\begin{frame}
\frametitle{Organized crime is not just limited to developing countries!} 
\begin{figure}[H]
  \includegraphics[keepaspectratio,width=0.8\textwidth]{fig1.png}
    \captionsetup{justification=raggedright,singlelinecheck=false}
\caption*{\footnotesize{Source: Global Organized Crime Index (https://ocindex.net/2023/heatmap/)}}
\end{figure}
\end{frame}


\begin{transitionframe}
\frametitle{Roadmap for today}
\tableofcontents
\end{transitionframe}

\section{Cause and consequences of individuals joining (organized) crime}
\subsection{Pinotti (2015): The Causes and Consequences of Organised Crime: Preliminary Evidence across Countries}
\begin{transitionframe}
\begin{center}
  \begin{huge}
    \textcolor{NTHU1}{%
  Pinotti (2015): \\ The Causes and Consequences of Organised Crime: Preliminary Evidence across Countries
    }
  \end{huge}
\end{center}
\end{transitionframe}

\begin{frame}
\frametitle{Q: What are the causes and consequences of organized crime?} 
Organized crime imposes serious consequences 
\begin{itemize}\setlength\itemsep{3pt}
\item Violent activities can destroy physical and human capital
\item Also increases uncertainties and riskiness $\to$ hurt long run investment and growth
\item Yet, relationship among organized crime, economics, and institutions still uncertain
\end{itemize} \medskip
This paper provides cross-country empirical evidence on these relations
\begin{itemize}\setlength\itemsep{3pt}
\item Organized crime $\propto $ levels of per capita economic output $\Downarrow$
\begin{itemize}
\item Two-way relation!! 
\end{itemize}
\item Politicians are more corrupt and/or vulnerable to attacks by organized crime
\begin{itemize}
\item Suggest that organized crime$\to$ Weak institution
\end{itemize}
\item Organized crime presence is positively correlated with natural resource endowment
\begin{itemize}
\item But only in fragile states with weak legal capacity: Weak institution$\to$ organized crime
\end{itemize}
\end{itemize}
\end{frame}

\begin{frame}
\frametitle{Measuring organized criminal group presence} 
By nature, measuring organized criminal group presence is tricky
\begin{itemize}\setlength\itemsep{3pt}
\item Usually do not operate out in the open (and not always violent)
\item When they do, they can pressure law enforcement to prevent detection
\end{itemize} \medskip
How this paper approaches this issue
\begin{itemize}\setlength\itemsep{3pt}
\item Source: Executive Opinion Survey (World Economic Forum data) across 147 countries
\begin{itemize}\setlength\itemsep{3pt}
  \item `In your country, to what extent does organized crime impose costs on business?'
  \item Measured on 1 (least affected) to 7 (most affected) scale, averaged across `06-`13

\end{itemize}
\item Economic performance from Penn World Table 8.0
\item Political variables from World Bank Governance Indicators
\item OLS to derive descriptive relations 
\[
Y_i = \alpha + \beta\text{Organized crime}_i +\gamma X_i + e_i
\]
\end{itemize}
\end{frame}

\begin{frame}
\frametitle{Organized crime correlates well with homicide rate} 
\begin{figure}[H]
\includegraphics[keepaspectratio,width=0.7\textwidth]{p2015_f0.jpeg}
\end{figure}
\end{frame}

\begin{frame}
\frametitle{Organized crime negatively correlates with economic indicators} 
\begin{figure}[H]
\includegraphics[keepaspectratio,width=0.7\textwidth]{p2015_f1.jpeg}
\end{figure}
\end{frame}

\begin{frame}
\frametitle{Organized crime negatively correlates with economic indicators} 
\begin{figure}[H]
\includegraphics[keepaspectratio,width=0.9\textwidth]{p2015_t1.jpeg}
\end{figure}
\end{frame}

\begin{frame}
\frametitle{Organized crime negatively correlates with institutional quality} 
\begin{figure}[H]
\includegraphics[keepaspectratio,width=\textwidth]{p2015_f2.jpeg}
\end{figure}
\end{frame}

\begin{frame}
\frametitle{Organized crime negatively correlates with institutional quality} 
\begin{figure}[H]
\includegraphics[keepaspectratio,width=0.9\textwidth]{p2015_t2.jpeg}
\end{figure}
\end{frame}

\begin{frame}
\frametitle{Positively correlated with natural resource dependence in weak states} 
\begin{figure}[H]
\includegraphics[keepaspectratio,width=0.7\textwidth]{p2015_f3.jpeg}
\end{figure}
\end{frame}

\begin{frame}
\frametitle{Discussion and takeaways} 
These findings highlight descriptive relations, but not causal
\begin{itemize}\setlength\itemsep{3pt}
\item Organized crime is hard to detect (measurement error)
\item Organized crime is a result of economic and political conditions (reverse causality)
\end{itemize}\medskip
Many modern research on the topic tries to unearth causal relationship
\begin{itemize}\setlength\itemsep{3pt}
\item Leverage sudden economic shock: Economic condition $\to$ Organized crime
\item Leverage political events (e.g. assassinations): Political event $\to$ Organized crime
\item IV method (e.g. road networks, agriculture): Organized crime $\to$ Economics and politics
\item In a rare instance, some people design an randomized experiments 
\end{itemize}\medskip
\end{frame}

\subsection{Chimeli, Soares (2017): The Use of Violence in illegal Markets: Evidence from Mahogany Trade in Brazilian Amazon}
\begin{transitionframe}
\begin{center}
  \begin{huge}
    \textcolor{NTHU1}{%
  Chimeli, Soares (2017): \\ The Use of Violence in illegal Markets: Evidence from Mahogany Trade in Brazilian Amazon
    }
  \end{huge}
\end{center}
\end{transitionframe}

\begin{frame}
\frametitle{Q: How does illegalization of a market cause violence?}
In illegal markets, state presence is limited 
\begin{itemize}\setlength\itemsep{3pt}
\item No formal means to enforce contracts and guarantee property rights
\item Violence potentially is used as arbitration device
\item But causal link between market illegality and violence is limited
\begin{itemize}
\item Extremely rare, unable to do randomized experiment
\end{itemize}
\end{itemize} \medskip
This paper uses illegalization of mahogany trade in Brazil to answer this question
\begin{itemize}\setlength\itemsep{3pt}
\item Leverage variation in homicide rates across areas with different mahogany shares
\item Homicide rates $\Uparrow$ in high mahogany growth areas 
\item This is amplified in states with high mahogany export shares
\item Homicides subsided with federal effort to reduce illegal mahogany trades itself
\end{itemize}
\end{frame}

\begin{frame}
\frametitle{Mahogany trades in Brazil and violence}
Brazil as a  Mahogany exporter
\begin{itemize}\setlength\itemsep{3pt}
\item One of the largest exporters before prohibition in 2001
\item Illegalization: Limited issuing of licenses in 1999, and introduced export quota at 2001
\item Illegalization of mahogany markets since
\begin{itemize}\setlength\itemsep{3pt}
\item Traded as `other timber species'
\item Trade picked up until active monitoring and enforcement came in at late in 2000s
\end{itemize}
 \end{itemize} \medskip
 Violence and illegal mahogany markets
 \begin{itemize}\setlength\itemsep{3pt}
\item No formal credit markets or contracts available $\to$ resort to violence
\item Use intimidation and threats on indigenous tribe leaders to steal Mahogany
\item Also use illegal weapons in the process
\item These were sometimes carried out by organized criminal groups in the area
 \end{itemize}
\end{frame}

\begin{frame}
\frametitle{Mahogany trades persisted under `others'}
\begin{figure}
  \includegraphics[keepaspectratio,width=0.75\textwidth]{cs2017_f1.jpeg}
\end{figure}
\end{frame}

\begin{frame}
\frametitle{Data and empirical strategy}
Data sources 
\begin{itemize}\setlength\itemsep{3pt}
\item  Mahogany presence: Brazilian Secretariat on International Trade
\item Violence: Ministry of Health keeps homicide rates at municipality level
\item Also control for mortality, land conflicts, total areas planted, agricultural GDP etc
\end{itemize} \medskip
Diff-in-diff leverages variation in mahogany and violence across time and municipalities
\[
Y_{it}= \alpha + \beta_1 D_{[`99\leq t\leq`01]\times}\text{M}_i + \beta_2 D_{[`02\leq t\leq`08]\times}\text{M}_i + \beta_3 D_{[t\geq`09]\times}\text{M}_i + \gamma Z_{it} + \phi_i+\mu_{st}+e_{it}
\]
\begin{itemize}\setlength\itemsep{3pt}
\item Breaks down into 3 periods: Different modes of illegalization and enforcement
\item Outcome variables: Primarily homicide rates
\end{itemize} \medskip
\end{frame}

\begin{frame}
\frametitle{Illegalization increases violence}
\begin{figure}
  \includegraphics[keepaspectratio,width=0.75\textwidth]{cs2017_t1.jpeg}
\end{figure}
\end{frame}

\begin{frame}
\frametitle{Particularly in a state with very high Mahogany exports}
\begin{figure}
  \includegraphics[keepaspectratio,width=0.7\textwidth]{cs2017_t2.jpeg}
\end{figure}
\end{frame}

\begin{frame}
\frametitle{Victims usually prime aged male}
\begin{figure}
  \includegraphics[keepaspectratio,width=0.75\textwidth]{cs2017_t3.jpeg}
\end{figure}
\begin{itemize}
\item Suggests that deaths are due to criminal dispute over domestic violence
\end{itemize}
\end{frame}

\begin{frame}
\frametitle{Discussion and takeaways}
Illegality of markets can increase violence
\begin{itemize}\setlength\itemsep{3pt}
\item These are not explained by socioeconomic conditions
\item Monitoring and enforcement $\to$ Illegal market $\Downarrow$ $\to$ Violence $\Downarrow$
\end{itemize} \medskip
What cautionary tale does it provide?
\begin{itemize}\setlength\itemsep{3pt}
\item In absence of adequate monitoring and enforcement, addressing unwanted externalities with overly restrictive regulations may increase social loss
\begin{itemize}\setlength\itemsep{3pt}
\item When markets become illegal, violence replaces formal credit markets and contracts
\end{itemize}
\item But how to establish capable monitoring? 
\begin{itemize}\setlength\itemsep{3pt}
\item Corruption can make monitoring difficult (police implicitly cooperating with bandits)
\end{itemize}
\end{itemize}
\end{frame}

\subsection{Dix-Carneiro, Soares, Ulyssea (2018): Economic Shocks and Crime: Evidence from Brazilian Trade Liberalization [Student presentation]}
\begin{transitionframe}
\begin{center}
  \begin{huge}
    \textcolor{NTHU1}{%
  Dix-Carneiro, Soares, Ulyssea (2018):\\ Economic Shocks and Crime: Evidence from Brazilian Trade Liberalization [Student presentation]
    }
  \end{huge}
\end{center}
\end{transitionframe}





\section{Criminal governance: Relationship between formal state and organized criminal groups}
\subsection{Dell (2015): Trafficking Networks and the Mexican Drug War}
\begin{transitionframe}
\begin{center}
  \begin{huge}
    \textcolor{NTHU1}{%
Dell (2015): \\ Trafficking Networks and the \\ Mexican Drug War
    }
  \end{huge}
\end{center}
\end{transitionframe}

\begin{frame}
\frametitle{Q: How do organized criminal groups respond to crackdowns?} 
Many countries engage in counterinsurgency strategies to drive out organized crime
\begin{itemize}\setlength\itemsep{3pt}
\item Mexico: Spent \$9 billion USD per year ($\simeq$ Social development expenses)
\item Causal evidence of crackdowns is rare
\begin{itemize}\setlength\itemsep{3pt}
\item Selection bias: State chooses where crackdown takes place deliberately
\item Reverse causality: Governments choose places likely to expect high violence
\end{itemize}
\end{itemize}\medskip
This paper uses plausibly exogenous variation in crackdowns 
\begin{itemize}\setlength\itemsep{3pt}
  \item Closely elected mayors of the party spearheading crackdowns (PAN)
  \item Places with close elections are less likely to be different on other dimensions
  \item Driven by drug traffickers killing each other (incumbent gang vs new gangs)
  \item PAN victories divert trafficking route, increase violence on alternate routes
\end{itemize}\medskip
\end{frame}

\begin{frame}
\frametitle{Background information} 
Drug Trafficking in Mexico
\begin{itemize}\setlength\itemsep{3pt}
  \item Large industry: \$14-48 billion USD (mostly to US)
  \item Drug trafficking organizations (DTOs) present in 68\% of Mexican municipalities
  \item Number of operating DTOs increasing since 2006 (also engage in other illicit activities)
\end{itemize}\medskip
Crackdowns on the DTOs
\begin{itemize}\setlength\itemsep{3pt}
  \item Felipe Calderón (PAN): Made crackdowns a centerpiece of his administration
 \begin{itemize}\setlength\itemsep{3pt}
  \item Mexico was previously a one-party (PRI) country, which took passive stance on DTOs
  \item Federal police/military arrest high-rank traffickers and seize drugs 
  \end{itemize}
\item Mayors: Name police officers who can assist federal government crackdowns
\begin{itemize}\setlength\itemsep{3pt}
  \item PAN mayors were actively supporting this
  \item Also more likely to request federal help compared to non-PAN mayors
\end{itemize}  
\end{itemize}\medskip
\end{frame}

\begin{frame}
\frametitle{Data and empirical strategy} 
Data sources
\begin{itemize}\setlength\itemsep{3pt}
  \item Official govt data on drug trades, DTO presence and related crimes (confidential)
  \item Election data from state electoral committees (available right now)
\end{itemize}\medskip
Effects of close PAN victory (i.e. increased crackdown)
\begin{itemize}\setlength\itemsep{3pt}
  \item Regression discontinuity leveraging close margin of victory for PAN
  \[
  Y_m = \alpha_0 + \alpha_1\text{PAN}_m+\alpha_2\text{PAN}_m\times f(\text{margin}_m)+\alpha_3(1-\text{PAN}_m)\times f(\text{margin}_m) +e_m
  \]
  \item Close election of PAN mayors vs Close loss of PAN mayors
  \item Other characteristics equal: Confirms that no other characteristics differ (Table 1)
\end{itemize}\medskip
\end{frame}



\begin{frame}
\frametitle{Existing drug-related violence intensifies} 
\begin{figure}
\includegraphics[keepaspectratio, width=\textwidth]{d2015_f1.jpeg}
\end{figure}
\begin{itemize}\setlength\itemsep{3pt}
\item Driven by rival gangs attacking incumbent gangs weakened by enforcements
\item Does not happen pre-election or pre-inauguration after elections (Rest of Fig 3)
\end{itemize}
\end{frame}

\begin{frame}
\frametitle{Do DTO responses have spillovers?} 
Logic behind spillovers to other regions
\begin{columns}
\begin{column}{0.49\textwidth}
\begin{figure}
\includegraphics[keepaspectratio,width=0.75\textwidth]{d2015_f2.jpeg}
\end{figure}  
\end{column}
\begin{column}{0.49\textwidth}
  \begin{itemize}
\item Close PAN victories raise costs of trafficking 
\item This incentivizes DTOS to shift paths 
\item Leading to increase in drug-related violence elsewhere  
  \end{itemize}
\end{column}
\end{columns}
\end{frame}

\begin{frame}
\frametitle{Indeed, violence do spillover} 
\begin{figure}
\includegraphics[keepaspectratio,width=0.7\textwidth]{d2015_t1.jpeg}
\end{figure} 
\end{frame}

\begin{frame}
\frametitle{Takeaways and discussions} 
Crackdowns induce direct and spillover costs 
\begin{itemize}\setlength\itemsep{3pt}
\item Violence in PAN municipalities $\Uparrow$ ($\because$ rival gangs kill weakened incumbent gangs)
\item Violence spills over to nearby municipalities ($\because$ reroute to find cheaper transport costs)
\item Thus, DTOs respond to changing environments, offsetting effectiveness of crackdowns
\end{itemize}\medskip
So what does it inform about crackdown policies?
\begin{itemize}\setlength\itemsep{3pt}
\item Emphasize ways to deter violence at its roots
\begin{itemize}
\item Other ways to disincentivize violence (very tricky question)
\end{itemize}
\item Focusing on high-rank arrests may open other gangs to take over
\begin{itemize}\setlength\itemsep{3pt}
\item Especially in places where inter-cartel competition is fierce (Mexico)
\item End up having criminal governance by rival gangs
\end{itemize}
\end{itemize} 
\end{frame}

\subsection{Blattman, Duncan, Lessing, Tobón (2025): Gang Rule: Understanding and Countering Criminal Governance}
\begin{transitionframe}
\begin{center}
  \begin{huge}
    \textcolor{NTHU1}{%
  Blattman, Duncan, Lessing, Tobón (2025): \\ : Gang Rule: Understanding and Countering Criminal Governance
    }
  \end{huge}
\end{center}
\end{transitionframe}

\begin{frame}
\frametitle{Q: Can state and criminal governance coexist? How?} 
Classical explanation focuses on the \textit{monopolitic} governance of state or criminal groups
\begin{itemize}\setlength\itemsep{3pt}
\item Fill in the vacuum: Criminal groups fill in the absence of the state
\item Competition: State exerts influence by driving out criminal groups (and vice versa)
\item But there are cases where \textit{duopolistic} rule exists (Italy, South Africa, Latam etc.)
\begin{itemize}
\item Criminal groups rule civilians and coexist with state governance
\end{itemize}
\end{itemize}\medskip
This paper identifies how duopolistic governance becomes possible
\begin{itemize}\setlength\itemsep{3pt}
  \item Gangs govern civilians to protect illegal revenues (drug dealing)
  \item In settings where such gains are high, state presence $\Uparrow$ $\Rightarrow$ Criminal rule $\Uparrow$
  \item Thus, state and criminal rule becomes complements, not substitutes
  \item Leverages exogenous boundary changes that changes the distance to closest state HQ
  \item Distance to HQ $\Downarrow$ $\to$ state responsiveness and criminal presence $\Uparrow$
\end{itemize}
\end{frame}

\begin{frame}
\frametitle{Context and framework} 
Context: Medellín, Colombia
\begin{itemize}\setlength\itemsep{3pt}
  \item Has well-organized bureaucracy and large fiscal revenues
  \item Security and legal services provided at sub-city level
  \item Poorer blocks have \textit{combos}: small local gangs providing security and drug trade
\end{itemize} \medskip
Why do gangs provide security services? (interview with gang leaders)
\begin{itemize}\setlength\itemsep{3pt}
  \item Providing security reduce need for police presence
  \item Expected residents to be loyal (i.e. cooperate with gangs over state)
\end{itemize}
How can state-gangs be complementary?
\begin{itemize}\setlength\itemsep{3pt}
  \item Classic: When state fails to meet demand for governance, gangs enter
  \item Twist: Simultaneous governance protects gangs illicit business from state repression
  \begin{itemize}\setlength\itemsep{3pt}
  \item Illicit profit $\Uparrow$ \& gang-state services differ: Gang incentive to govern areas close to state $\Uparrow$
  \item Low profit or state already dominant: Incentive to govern $\Downarrow$
  \end{itemize}
\end{itemize}
\end{frame}

\begin{frame}
\frametitle{Data and empirical strategy} 
Data sources
\begin{itemize}\setlength\itemsep{3pt}
\item Interviews with the criminal group leaders to understand motives and combo presence
\item Surveys on residents and businesses to gather data on governance and perceptions
\end{itemize}\medskip
Geographic regression discontinuity that leverages 1987 reform 
\begin{itemize}\setlength\itemsep{3pt}
  \item Analysis at block-pairs
  \item New borders in Medellín (6 sub-area $\to$ 16) divided a subregion into 2-3 smaller units
  \item Changes the nearest security and justice service location for some blocks
  \item Regress the following equation: $\Delta d_{ij}$ is the treatment  ($i,j$: block pairs, $b$: border)
  \[
  \Delta Y_{ijb} = \alpha_b + \beta \Delta d_{ij}+\theta \Delta X_{ij}+ \lambda B_{ij}+e_{ijb}
  \]
  \item Confirm that other characteristics do not change with the new border introduction
\end{itemize}
\end{frame}

\begin{frame}
\frametitle{Description of the regression discontinuity} 
\begin{figure}[H]
\includegraphics[keepaspectratio, width=\textwidth]{bdlt2025_f1.jpeg}
\end{figure}
\end{frame}

\begin{frame}
\frametitle{State and combo presence increases as distance to HQ decreases} 
\begin{figure}[H]
\includegraphics[keepaspectratio, width=0.6\textwidth]{bdlt2025_t1.jpeg}
\end{figure}
\end{frame}

\begin{frame}
\frametitle{Combo response strong in areas with high drug rents} 
Regression augmented with measure of illicit rent $\pi_{ij}$
\[
  \Delta Y_{ijb} = \alpha_b + \beta \Delta d_{ij} 
  +\gamma \pi_{ij}+ \delta \Delta d_{ij}\times \pi_{ij}+\theta \Delta X_{ij}+ \lambda B_{ij}+e_{ijb}
\]
\begin{figure}[H]
\includegraphics[keepaspectratio, width=\textwidth]{bdlt2025_t2.jpeg}
\end{figure}
\end{frame}

\begin{frame}
\frametitle{Presence of criminal groups does not increase perceived legitimacy} 
Legitimacy: combination of trust, fairness, and satisfaction
\begin{itemize}\setlength\itemsep{3pt}
\item State and combo legitimacy does not increase (two slides before)
\item In high drug rent areas, closer $\Rightarrow$ legitimacy for both state and combos $\Downarrow$ (slide before)
\item Qualitative interview results match: Residents unhappy with influx of `undesirable' ppl
\item No significant economic development for most measures...
\end{itemize}\medskip
Why did this happen?
\begin{itemize}\setlength\itemsep{3pt}
\item Gangs overestimating how loyal citizens were to them
\item Gangs entering areas to offset loyalty penalty coming from drug sales
\item Yet, providing security could reduce police patrol (so motives to enter exist)
\end{itemize}\medskip
\end{frame}

\begin{frame}
\frametitle{Takeaways and discussions} 
State and criminal presence can be complementary
\begin{itemize}\setlength\itemsep{3pt}
  \item Particularly in areas where returns from illegal trades are high
  \item Could explain converse case: Low illicit returns $\Rightarrow$ Gang extortions rise
\end{itemize}\medskip
What is to be learned from this? (ongoing question)
\begin{itemize}\setlength\itemsep{3pt}
\item To cut criminal presence, need to attack root causes that incentivize them
\begin{itemize}\setlength\itemsep{3pt}
\item Tougher than said, what are the criminal groups after to begin with?
\item Criminal groups can strategically move to retain these benefits (see any drug wars)
\end{itemize}
  \item Can the findings be generalized, even in places with high extent of illicit activities?
 \begin{itemize}
\item Active state vs. organized criminal conflict in such places too (see Mexico)
\end{itemize}
\end{itemize}\medskip
\end{frame}


\subsection{Trejo, Ley (2018): Why Did Drug Cartels Go to War in Mexico? Subnational Party Alternation, the Breakdown of Criminal Protection, and the Onset of Large-Scale Violence [student presentation]}
\begin{transitionframe}
\begin{center}
  \begin{huge}
    \textcolor{NTHU1}{%
Trejo, Ley (2018): \\ Why Did Drug Cartels Go to War in Mexico? Subnational Party Alternation, the Breakdown of Criminal Protection, and the Onset of Large-Scale Violence [student presentation]
    }
  \end{huge}
\end{center}
\end{transitionframe}

\begin{frame}
\frametitle{Takeaways and remaining questions} 
Many individuals join organized crime, some of which evolve to criminal governance
\begin{itemize}\setlength\itemsep{3pt}
\item Particularly in economically or institutionally less developed areas
\item Illegalization of the markets potentially gives organized crime an opportunity
\item Some of them are on verge to compete governments out
\begin{itemize}
\item Either strategically coexisting in and duopoly
\item Or adjusting their survival strategy accordingly
\item Sometimes outright fighting formal state
\end{itemize}
\end{itemize}\medskip
Unanswered questions
\begin{itemize}\setlength\itemsep{3pt}
\item How do you effectively curb incentives to join crime?
\begin{itemize}\setlength\itemsep{3pt}
\item Monitoring requires transparent and accountable government: Not always the case
\item Hardline stance can go wrong (Mexico) but it can deliver a strong message 
\end{itemize}
\item What can be done to address criminal governance?
\begin{itemize}\setlength\itemsep{3pt}
\item Organized crime not accountable to the populace: How to break up these groups?
\item Minimizing state absence/capture: Organized crime presence limited in capable states
\end{itemize}
\end{itemize}
\end{frame}


%%%%%%%%%%%
\end{document}


